\documentclass[12pt]{article}
\usepackage[utf8]{inputenc}
\usepackage[T1]{fontenc}
\usepackage{hyperref}
\title{Kariyarra Language Resources for K-12 Curriculum Development}
\author{Compiled by Claude Sonnet 5 (claude-sonnet-5) for Andr\'as Kornai}
\date{}
\begin{document}
\maketitle

\section{Introduction}
Kariyarra (ISO 639-3 vka) is a Pama-Nyungan language of the Port Hedland area in the Pilbara region of Western Australia. It is closely related to Ngarluma and is regarded by some linguists as a dialect in a continuum with it. It is highly endangered or dormant. One search result estimated fewer than two fluent speakers in 2002. The Town of Port Hedland states that Kariyarra is the language of its region.

\section{Primary Educational Resources}
No school programme or textbook was found. The Wangka Maya materials below are the closest educational resource.

\section{Linguistic and Reference Materials}
\begin{itemize}
\item Wangka Maya Pilbara Aboriginal Language Centre, Kariyarra page, with written example sentences: \url{https://www.wangkamaya.org.au/pilbara-languages/kariyarra}
\item Wangka Maya, Kariyarra cultural introduction: \url{https://www.wangkamaya.org.au/pilbara-languages/kariyarra/01-cultural-introduction}
\item Wikipedia, Ngarluma language (for the related dialect): \url{https://en.wikipedia.org/wiki/Ngarluma_language}
\end{itemize}

\section{Religious and Community Resources}
\begin{itemize}
\item Town of Port Hedland, Aboriginal Beginnings: \url{https://www.porthedland.wa.gov.au/our-history.aspx}
\item Town of Port Hedland, Ngalikuru Kuma Marri Draft Plan 2024--2028, which acknowledges Kariyarra, Ngarla and Nyamal Traditional Owners: \url{https://www.porthedland.wa.gov.au/Profiles/porthedland/Assets/ClientData/Ngalikuru_Kuma_Marri_Draft_Plan_2024-2028.pdf}
\end{itemize}
The Kariyarra Aboriginal Corporation holds native title for the Kariyarra people. A website for it was not found in the search results, so none is cited.

\section{Digital and Audio Resources}
The Wangka Maya Kariyarra pages are described as including recordings. The search did not confirm separate audio URLs.

\section{Curriculum Development Context}
Kariyarra has very few resources, so a heritage curriculum would be a revitalization effort led by the Kariyarra community, with Wangka Maya as the natural partner. It would draw on the existing Wangka Maya materials and on Ngarluma materials where the community considers them appropriate.
\end{document}